\documentclass[11pt]{article}

\usepackage[margin=1in]{geometry}
\usepackage{booktabs}
\usepackage{array}
\usepackage[hidelinks]{hyperref}

\newcolumntype{L}[1]{>{\raggedright\arraybackslash}p{#1}}
\emergencystretch=3em

\title{Supplementary Material for IAD-Risk}
\author{Anonymous Authors}
\date{}

\begin{document}
\maketitle

\section{Scope}

This supplementary material documents the reproducibility boundary for the manuscript ``IAD-Risk: Risk-Aware Identity-Agenda Disentanglement for Scholarly Work Deduplication.'' It is intended to accompany the main manuscript, not to replace the project repository or the final artifact release. The repository contains executable code, fixture data, schema contracts, and build scripts; large third-party raw data, model checkpoints, and full experimental outputs are kept outside Git and must be distributed through an artifact release when numerical result reproduction is required.

\section{Reproduction Levels}

\begin{table}[h]
\caption{Reproduction levels and expected evidence.}
\centering
\small
\begin{tabular}{L{0.12\linewidth}L{0.35\linewidth}L{0.43\linewidth}}
\toprule
Level & Required input & Evidence produced \\
\midrule
L0 & Source tree and configured Python environment & CLI help, public-release scan, unit tests, and PDF build logs. \\
L1 & Public miniature fixtures in \texttt{tests/fixtures} & Eval-source files and an IAD-Bench fixture package with the documented schema. \\
L2 & Independently obtained public raw files & Rebuilt DeepMatcher, SciRepEval-style, OpenAlex-style, and OpenCitations-derived evaluation sources. \\
L3 & Released artifact package & Fixed tables, predictions, logs, manifests, checksums, and commit identifiers for result audit. \\
\bottomrule
\end{tabular}
\end{table}

\section{No-Network Fixture Rebuild}

The fixture rebuild validates the data adapters and IAD-Bench schema without depending on external network access. It is not a substitute for the full paper experiment.

\small
\begin{verbatim}
python -m iad_sieve.cli prepare-deepmatcher \
  --table-a tests/fixtures/deepmatcher/tableA.csv \
  --table-b tests/fixtures/deepmatcher/tableB.csv \
  --pairs tests/fixtures/deepmatcher/test.csv \
  --dataset-name deepmatcher_fixture \
  --output-dir outputs/repro_fixture/deepmatcher

python -m iad_sieve.cli prepare-scirepeval-proximity \
  --metadata tests/fixtures/scirepeval/metadata.jsonl \
  --pairs tests/fixtures/scirepeval/scidocs_cite_pairs.csv \
  --dataset-name scirepeval_fixture \
  --output-dir outputs/repro_fixture/scirepeval

python -m iad_sieve.cli prepare-openalex-weak-labels \
  --works tests/fixtures/openalex/works.jsonl \
  --citations tests/fixtures/openalex/coci.csv \
  --dataset-name openalex_fixture \
  --output-dir outputs/repro_fixture/openalex \
  --min-shared-references 1 \
  --max-pairs 20

python -m iad_sieve.cli build-iad-bench \
  --source-dirs \
    outputs/repro_fixture/deepmatcher \
    outputs/repro_fixture/scirepeval \
    outputs/repro_fixture/openalex \
  --output-dir outputs/repro_fixture/iad_bench \
  --seed 42
\end{verbatim}
\normalsize

The manuscript package includes a verification wrapper for the same no-network path:

\small
\begin{verbatim}
python manuscript/scripts/verify_fixture_rebuild.py
\end{verbatim}
\normalsize

The wrapper runs the four fixture commands in a temporary directory and checks that each source adapter writes \texttt{eval\_documents.jsonl}, \texttt{eval\_pairs.jsonl}, and \texttt{dataset\_summary.jsonl}, and that the assembled IAD-Bench package writes pair, document, split, summary, provenance, and dataset-card files.

\section{Public-Source Rebuild}

For full data reconstruction, raw public files should be obtained from their original sources and placed under local \texttt{data/raw/} directories. These directories are intentionally excluded from Git. The project keeps conversion commands and schema contracts under version control, while source licenses and access conditions remain governed by the original providers.

\small
\begin{verbatim}
python -m iad_sieve.cli prepare-deepmatcher \
  --table-a data/raw/deepmatcher/DBLP-ACM/tableA.csv \
  --table-b data/raw/deepmatcher/DBLP-ACM/tableB.csv \
  --pairs data/raw/deepmatcher/DBLP-ACM/test.csv \
  --dataset-name DBLP-ACM \
  --output-dir data/processed/eval_sources/deepmatcher_dblp_acm

python -m iad_sieve.cli prepare-scirepeval-proximity \
  --metadata data/raw/scirepeval/metadata.jsonl \
  --pairs data/raw/scirepeval/scidocs_cite_pairs.csv \
  --dataset-name scidocs_cite \
  --output-dir data/processed/eval_sources/scirepeval_scidocs_cite

python -m iad_sieve.cli prepare-openalex-weak-labels \
  --works data/raw/openalex/works_cs_sample.jsonl \
  --citations data/raw/opencitations/coci.csv \
  --dataset-name openalex_cs_sample \
  --output-dir data/processed/eval_sources/openalex_cs_sample \
  --min-shared-references 1 \
  --max-pairs 5000

python -m iad_sieve.cli build-iad-bench \
  --source-dirs \
    data/processed/eval_sources/deepmatcher_dblp_acm \
    data/processed/eval_sources/scirepeval_scidocs_cite \
    data/processed/eval_sources/openalex_cs_sample \
  --output-dir data/processed/iad_bench/open_v3 \
  --train-ratio 0.8 \
  --dev-ratio 0.1 \
  --seed 42
\end{verbatim}
\normalsize

\section{Artifact Package Requirements}

Full numerical result reproduction should use a released artifact package rather than raw Git-tracked files. A complete artifact package should contain a manifest, checksums, tables, predictions, configuration files, logs, and the commit identifier used to generate the results.

\small
\begin{verbatim}
paper-artifacts/
  README.md
  manifest.json
  checksums.sha256
  configs/
  tables/
  predictions/
  reports/
  logs/
\end{verbatim}
\normalsize

The manifest should record the source data version, processing commands, dependency versions, random seeds, thresholds, split configuration, hardware summary, and SHA256 checksum for every result file. The checksums should be validated before using artifact files to support claims in the manuscript.

\section{Claim-Evidence Matrix}

Table~\ref{tab:claim-evidence} maps the major manuscript claims to the evidence required for responsible interpretation. The matrix is part of the submission boundary: a claim should be weakened or removed if the corresponding evidence file is absent from the released artifact package.

\begin{table}[h]
\caption{Claim-evidence matrix for the main manuscript.}
\label{tab:claim-evidence}
\centering
\small
\begin{tabular}{L{0.26\linewidth}L{0.30\linewidth}L{0.34\linewidth}}
\toprule
Claim & Required evidence & Boundary \\
\midrule
Identity-agenda confusion is a practical false-merge risk. & Open-v2 hard-negative evaluation with FMR and HNFMR fields. & Supports a targeted scholarly deduplication failure mode, not a prevalence estimate for all domains. \\
IAD-Risk suppresses hard-negative false merges under the reported setting. & IAD-Risk prediction files, metric summaries, split identifiers, thresholds, and checksums. & Supports targeted false-merge suppression; does not establish broad method superiority. \\
IAD-Bench is provenance-aware. & Pair schema, label strength, label source, provenance fields, and fixture rebuild commands. & Supports a benchmark contract; does not turn silver or proxy labels into human gold. \\
Repository-level reproduction is possible without committing raw data. & L0/L1 fixture rebuild, public-release scan, unit tests, and documented public-source commands. & Verifies code and schema behavior; full numerical reproduction requires L2/L3 source files or artifacts. \\
\bottomrule
\end{tabular}
\end{table}

\section{Uncertainty and Ablation Requirements}

Result packages that report confidence intervals should include bootstrap outputs generated from the exact prediction files used in the manuscript tables. The repository provides \emph{run-iad-evidence-bootstrap} for stratified pair-level intervals and \emph{run-bootstrap} for general metric intervals. A numerical interval should be cited only when the corresponding CSV file, command log, random seed, and checksum are present in the artifact package.

Mechanism claims should be tied to explicit ablation outputs. The repository provides \emph{run-iad-ablation-suite} for comparing full IAD-Risk decisions with variants that remove risk gating, agenda-non-identity blocking, or single-space separation. Without those outputs, the main manuscript should keep the mechanism claim at the level of stratified HNFMR evidence rather than claiming a complete ablation study.

\section{Claim Boundary}

The main manuscript treats DeepMatcher-style labels as gold identity evidence. It treats SciRepEval and SciDocs-style proximity as same-agenda proxy evidence. It treats OpenAlex and OpenCitations-derived pairs as silver hard-negative evidence. OpenAlex and OpenCitations signals are not human gold labels. Broad method-ranking claims require stronger source-heldout evidence and manual validation beyond the current manuscript package.

\end{document}
